\documentclass[letterpaper,10pt,onecolumn]{article}
\usepackage[spanish]{babel}
\usepackage[latin1]{inputenc}
\usepackage[pdftex]{color,graphicx}
\usepackage{hyperref}
\setlength{\oddsidemargin}{0cm}
\setlength{\textwidth}{490pt}
\setlength{\topmargin}{-40pt}
\addtolength{\hoffset}{-0.3cm}
\addtolength{\textheight}{4cm}

\begin{document}
\begin{center}

\includegraphics[width=490pt]{header.png}\\[0.5cm]

\textsc{\LARGE Astronom\'ia Popular}\\[0.1cm]

\large Jaime E. Forero Romero\\[0.5cm]

\end{center}

\large \noindent\textsc{Nombre del curso:} Astronom\'ia Popular %Aqui
                                %nombre del curso 
 
\noindent\textsc{C\'odigo del curso:} FISI-XXXX. CBU Tipo E (intensivo en
escritura acad\'emica).  Como tipo A/B dependiendo de la
clasificaci\'on del comit\'e de CBU.  %Aqui el 
                                %codigo del curso 

\noindent\textsc{Unidad acad\'emica:} Departamento de F\'isica 

\noindent\textsc{Periodo acad\'emico:} 201520 %Aqui el periodo,
                                %p.ej. 201510 

\noindent\textsc{Horario:} %Aqui el horario, p.ej. Ma y Ju, 10:00 a
                           %11:20 

\noindent\rule{\textwidth}{1pt}\\[-0.3cm]

\normalsize \noindent\textsc{Nombre profesor(a) principal:} Jaime
E. Forero-Romero %Aqui nombre del profesor principal 

\noindent\textsc{Correo electr\'onico:}
\href{mailto:je.forero@uniandes.edu.co}{\nolinkurl{je.forero@uniandes.edu.co}}
%Cambie address por su direccion de correo uniandes 

\noindent\textsc{Horario y lugar de atenci\'on:} Ip208 %Aqui su
                                %horario y lugar de atencion,
                                %p.ej. Vi, 15:00 a 17:00, Oficina
                                %Ip102 
\\[-0.1cm]


\noindent\rule{\textwidth}{1pt}\\[-0.1cm]

\newcounter{mysection}
\addtocounter{mysection}{1}

\noindent\textbf{\large \Roman{mysection} \quad
  Introducci\'on}\\[-0.2cm] 

%Este espacio es para hacer una introduccion al curso, evidenciando la propuesta metodologica. Debe ser clara y precisa.

\noindent\normalsize 
La astronom\'ia est\'a presente en la informaci\'on que
recibimos a diario en el Siglo XXI. Posible vida en otros planetas, el
origen del  Universo, apocalipsis astron\'omicos, ciencia ficci\'on,
celebraciones rituales en el calendario y la exploraci\'on espacial
son algunos de los temas que parecen familiares a la mayor\'ia pero que
en realidad no los conocemos en toda la profundidad que merecen. El
objetivo principal del curso es tratar aspectos de la cultura popular
relacionados con la astronom\'ia y presentarlos a trav\'es del
conocimiento que se ha constru\'ido en cosmolog\'ia, astronom\'ia
gal\'actica, astrof\'isica planetaria, astrof\'isica estelar,
astronom\'ia de posici\'on, coheter\'ia, ingenier\'ia aeroespacial y
astrobiolog\'ia. \\[0.1cm]  

\stepcounter{mysection}
\noindent\textbf{\large \Roman{mysection} \quad Objetivos}\\[-0.2cm]

%En este espacio se debe precisar el ente visor del curso y el proposito ideal al finalizar el curso.

\noindent\normalsize Los objetivos principales del curso son:

\begin{itemize}
\item Reconocer el lugar del conocimiento astron\'omico en diferentes
  aspectos de la cultura popular. \\[-0.6cm] 

\item Examinar los v\'inculos hist\'oricos de la astronom\'ia con
  otras \'areas del conocimiento.  \\[-0.6cm] 

\item Reconocer diferentes maneras de construir y validar conocimiento
  en las ciencias naturales, en particular en astronom\'ia.\\[-0.2cm]
\end{itemize}

\stepcounter{mysection}
\noindent\textbf{\large \Roman{mysection} \quad Competencias a
  desarrollar}\\[-0.2cm] 

%En este espacio se describen las habilidades que el estudiante
%desarrollara en el transcurso del curso. 

\noindent\normalsize Al finalizar el curso, se espera que el
estudiante est\'e en capacidad de: 

\begin{itemize}
\item Analizar contenidos de diferentes fuentes en la cultura popular
  para encontrar su relaci\'on con el conocimiento astron\'omico.\\[-0.6cm]
\item Leer textos no especializados con contenido astron\'omico para
  reconocer el posible uso correcto o incorrecto de conceptos cient\'ificos.
\\[-0.6cm]
\item Escribir un texto sobre alg\'un aspecto de la cultura popular
  influenciado por la astronom\'ia para evidenciar la relaci\'on del
  conocimiento cient\'ifico con las creencias/valores del p\'ublico no
  especializado.\\[-0.2cm]  
\end{itemize}

\stepcounter{mysection}
\noindent\textbf{\large \Roman{mysection} \quad Contenido por
  semanas}\\[-0.2cm] 

%Se expone de forma ordenada toda la tematica a tratar del curso. Debe
%planearse para 15 semanas. 

\noindent\normalsize \textbf{\textsc{Semana 1.}} El Comienzo. Las
diferentes historias de origen en los \'ultimos dos mil a\~nos
en diferentes civilizaciones de Asia, Africa y Am\'erica. La teor\'ia
del Big Bang.\\[-0.3cm]   


\noindent\textbf{\textsc{Semana 2.}} Espacio y tiempo. Modelos de
espacio y tiempo desde Grecia antigua hasta Einstein. \\[-0.3cm] 

\noindent\textbf{\textsc{Semana 3.}} Ritmos de la vida diaria. Ciclos
lunares y solares que dictan ritmos biol\'ogicos y
dan estructura a sociedades. \\[-0.3cm] 

\noindent\textbf{\textsc{Semana 4.}} Ritmos de la vida
religiosa. Bases astron\'omicas presentes en la celebraci\'on de ritos
religiosos. \\[-0.3cm] 


\noindent\textbf{\textsc{Semana 5.}} Signos. Lo que la ciencia
contempor\'anea tiene que decir sobre la astrolog\'ia. \\[-0.3cm] 

\noindent\textbf{\textsc{Semana 6.}} El Espacio visto por los
pol\'iticos. La carrera espacial y su lugar en luchas de poder global
durante el Siglo XX y XXI.\\[-0.3cm]  

\noindent\textbf{\textsc{Semana 7.}} El Universo visto por los
periodistas. Los errores y aciertos de los periodistas en
cubrimientos de noticias sobre astronom\'ia. \\[-0.3cm]   

\noindent\textbf{\textsc{Semana 8.}} Vida m\'as all\'a de la
Tierra. Astrobiolog\'ia b\'asica y la b\'usqueda cient\'ifica de vida
extraterrestre. \\[-0.3cm]  

\noindent\textbf{\textsc{Semana 9.}} Nuevos mundos. Exoplanetas y la
b\'usqueda de nuevos mundos habitables .\\[-0.3cm] 

\noindent\textbf{\textsc{Semana 10.}} Humanos en el espacio. Historias
de cosmonautas y astronautas. La posible colonizaci\'on del sistema
solar. \\[-0.3cm] 

\noindent\textbf{\textsc{Semana 11.}} El espacio exterior en la
ciencia ficci\'on contempor\'anea. Paseo por la obra de Isaac Asimov,
Carl Sagan y Philip K. Dick.\\[-0.3cm] 

\noindent\textbf{\textsc{Semana 12.}} De pel\'icula. Referencias en el
cine de ciencia ficci\'on a la exploraci\'on espacial y viaje
transdimensional. \\[-0.3cm]  

\noindent\textbf{\textsc{Semana 13.}} Banda Sonora. Referencias de la
m\'usica popular a temas astron\'omicos. La ciencia detr\'as de
diferentes composiciones musicales. \\[-0.3cm] 

\noindent\textbf{\textsc{Semana 14.}} Space-Art. El punto de encuentro
entre la astronom\'ia y las ciencias del espacio con las artes
pl\'asticas.\\[-0.3cm] 

\noindent\textbf{\textsc{Semana 15.}} El Fin. Historias de apocalipsis
astron\'omicos y su verdadera probabilidad de suceder.\\[0.1cm] 

\stepcounter{mysection}
\noindent\textbf{\large \Roman{mysection} \quad Metodolog\'ia}\\[-0.2cm]

%Se describen las tecnicas y metodos para el desarrollo exitoso del curso.

\noindent\normalsize 
Cada semana los temas se presentar\'an por el instructor con
exposiciones orales y apoyo audiovisual (transparencias, audio,
video). En cada sesi\'on se presentar\'an dos aspectos del
conocimiento: el popular y el cient\'ifico. En cada tema se har\'a
\'enfasis en los aspectos que son apoyados o entran en contradicci\'on con el conocimiento cient\'ifico, as\'i como sobre
aquellas \'areas que se encuentran por fuera del alcance de la ciencia
contempor\'anea. El tono de las presentaciones ser\'a divulgativo,
presentando los conceptos importantes a trav\'es
de ejemplos y analog\'ias que no requieran un formalismo matem\'atico
complejo. Por su parte los estudiantes deben escribir tres ensayos
durante el curso. En tres sesiones del semestre se har\'an comentarios breves
(de media hora) sobre las entregas de los estudiantes. Por fuera de
clase se dar\'a una retroalimentaci\'on detallada por escrito a cada  ensayo.  
\\[0.1cm]

\stepcounter{mysection}
\noindent\textbf{\large \Roman{mysection} \quad Criterios de evaluaci\'on}\\[-0.2cm]

%Tener en cuenta los siguientes aspectos:
%	\item Porcentajes de cada evaluacion
%	\item Fechas importantes
%	\item Parametros de calificacion
%	\item Calificacion de asistencia y/o participacion en clase
%	\item Reclamos
%	\item Politica de aproximacion de notas

\noindent\normalsize 
El curso se calificar\'a con la entrega de tres ensayos de 2000 a 2500
palabras, cada uno de un valor del $33.3\%$ de la nota final.  Para
que los estudiantes tengan oportunidad de mejorar sus ensayos se
har\'an dos entregas. La primera servir\'a para dar una
retroalimentaci\'on inicial y tendr\'a una calificaci\'on de $10\%$
sobre la nota final. La segunda entrega ser\'a la definitiva y
tendr\'a  valor de $23.3\%$ sobre la nota final.  
 
El primer ensayo es un texto comparativo. Retoma un tema 
tratado en clase. El estudiante debe ampliar las referencias y la
discusi\'on hecha por el instructor.

El segundo ensayo es de car\'acter divulgativo. El estudiante retoma
un tema tratado en clase y lo convierte en un texto que podr\'ia
encontrarse en un medio de comunicaci\'on masivo.   

El tercer ensayo es de car\'acter argumentativo. El estudiante propone
su visi\'on sobre un tema nuevo que concuerda con el programa general
del curso. El ensayo debe tomar en cuenta una fuente de la cultura
popular (peri\'odicos, televisi\'on, radio, cine, c\'omics, museos) que
no haya sido explorada en profundidad durante el curso.   
\\[0.1cm]

\stepcounter{mysection}
\noindent\textbf{\large \Roman{mysection} \quad Bibliograf\'ia}\\[-0.2cm]

%Indicar los libros y la documentacion guia.


\noindent\normalsize Bibliograf\'ia principal:

\begin{itemize}
\item Carl Sagan. \textit{Cosmos: una evoluci�n c�smica de
  quince mil millones de a�os que ha transformado la materia
  en vida y en consciencia}, 1982. (Biblioteca General -
  523.1 S131 Z258 )
\end{itemize}

\noindent\normalsize Bibliograf\'ia complementaria:

\begin{itemize}
\item Roger D. Launius (Editor). \textit{Exploring the solar system:
  The history and science of planetary exploration},
  2013. (Biblioteca General - Recurso Electr\'onico 523.2 23 )\\[-0.6cm]
\item Bryan A. Penprase. \textit{The power of stars: how celestial
  observations have shaped civilization}, 2011. (Biblioteca General -
  Recurso Electr\'onico 523.1)\\[-0.6cm]
\item Carl Sagan. \textit{El mundo y sus demonios: la ciencia
  como una luz en la oscuridad}, Editorial Planeta,
  1997. (Biblioteca General - 001.94 S131 Z282) \\[-0.6cm]
\item Andrew M. Shaw. \textit{Astrochemistry: from astronomy to astrobiology}, 2011. (Biblioteca General -
  523.02 S318 )\\[-0.6cm]
\item Charles Percy Snow. \textit{The two cultures},
  1993. (Biblioteca General - 301.2 S558T 1993)\\[-0.6cm]
\end{itemize}

\end{document}
